\begin{abstract}
Memory-based embodied agents face one decision over and over: commit to an
object they remember, or keep exploring. We show that this decision, not what
is stored or retrieved, is where a strong memory agent fails. Labelling every
commitment of MemoryExplorer on LMEE-Bench with simulator ground truth, we find
that 57.4\% of navigation failures are commitments to the wrong remembered
object, three quarters of them made before the goal had been seen. On a
58-task subset, most such choices lie in a region other than the goal's.
Until the goal's region has been
observed, the agent commits as often when the goal is absent from memory as
when it is present. Keeping the whole history in the model's context, as
recent implicit memories do, makes matters worse (19.6 vs.\ 25.9 SR, same model)
unless the context is scoped to the goal's region: with ground-truth regions,
this adds 18.0 SR. We
therefore propose \ours, a region-scoped context memory. The VLM's context holds
one panorama per region, with unobserved space left visible as holes, and the
agent commits in two steps before it moves: first a region or ``explore'', then
a target inside that region only. \ours{} is trained on streams in which the
goal is often absent, and it needs no retrieval index, no stored features and
at most two VLM calls per decision. On LMEE-Bench and GOAT-Bench, \ours{}
improves success rate by \tbd{+X} over MemoryExplorer and cuts premature
commitments from \tbd{X}\% to \tbd{Y}\%.
\end{abstract}
